\chapter{Per-model atlas: the successful \texorpdfstring{$T=29$}{T=29} long-task models}
\label{app:atlas_t29}

This appendix gives the same per-model treatment as Appendix~\ref{app:atlas}, for the six
models that trained successfully on the \emph{long} ($T=29$, random-onset) change-detection
task: \code{v5}, \code{v5\_part2}, \code{v8}, \code{v11}, \code{v11\_part2}, and
\code{v11\_part5}. These predate the paper-faithful build and live in separate codebases;
they were each driven through one uniform analysis (the shared \code{\_behav\_utils} harness,
or \code{distractor\_dd} for \code{v11\_part5}) so the figures are directly comparable. Unlike
the four-patch paper network, these tokenize the $50\times50$ image into a $10\times10$ grid
($100$ tokens), so the attention maps are $10\times10$ and far richer.

\paragraph{Two architectural families.} The \emph{cross-attention} models (\code{v5},
\code{v5\_part2}, \code{v8}) place the image patches \emph{and} the memories in the key set,
so the attention has separate \emph{image-key} and \emph{memory-key} ($H_1,H_2$) maps --- both
are shown. The \emph{dual-stream} models (\code{v11}, \code{v11\_part2}, \code{v11\_part5})
use the image as the \emph{query} and attend over memory only, so their maps are per-stream
(salience$\to H_1/H_2$, top-down$\to H_2$): they show \emph{where in memory} the image query
looks. \code{v11\_part5} was trained on a cued-\emph{side} task with distractors; for
comparability it is evaluated here on the standard cued task with distractors switched off.

\paragraph{The headline, before the details.} On the long task these models \emph{do} orient
attention to the cue: the salience / image-key attention is biased toward the cued quadrant
\emph{before} the change (visible in panel~D of each behaviour figure), and the reaction time
is strongly graded with change magnitude (panel~C, $\sim$$6\to1$ frames) rather than floored.
This is the regime contrast anticipated by the short-task study (open question O5): the long
horizon both permits and rewards pre-orienting, where the seven-step task did not.

\begin{table}[h]
\centering
\caption{The six successful $T=29$ models. Mean cued-hit and false-alarm are averaged over
the validity$\times$magnitude psychometric grid (so the absolute hit is below the per-condition
peak, which approaches $1.0$ at large $\Delta$). Snapshot iteration as analysed.}
\label{tab:t29results}
\small
\begin{tabular}{llccc}
\toprule
Model & Family & snapshot iter & mean cued-hit & mean FA \\
\midrule
\code{v5}        & cross-attention (memory-as-tokens) & $1999$  & $0.750$ & $0.038$ \\
\code{v5\_part2} & cross-attention over memory        & $4999$  & $0.714$ & $0.044$ \\
\code{v8}        & cross-attention, $H_1$ residual    & $4999$  & $0.721$ & $0.019$ \\
\code{v11}       & parallel dual-stream               & $8199$  & $0.671$ & $0.059$ \\
\code{v11\_part2}& split-readout cross-talk           & $10399$ & $0.683$ & $0.025$ \\
\code{v11\_part5}& split-readout (distractor task)    & $11799$ & $0.748$ & $0.022$ \\
\bottomrule
\end{tabular}
\end{table}

\begin{center}\archlegend\end{center}

% \atlast {file-stem}{Title}{note}
\newcommand{\atlast}[3]{%
  \section{#2}%
  #3\par\medskip
  \begin{figure}[h]\centering\resizebox{\textwidth}{!}{\input{figures/arch_t29_#1}}%
    \caption{\textbf{#2} --- architecture. Conv patch-embed $\to$ $10\times10$ tokens; solid
    arrows forward, dashed red recurrent feedback.}\end{figure}
  \begin{figure}[h]\centering\includegraphics[width=\textwidth]{atlas_t29_#1_attention.pdf}%
    \caption{\textbf{#2} --- attention maps at every frame ($T=29$), $100\%$ cue at $S_1$, for
    the no-change and change-at-$S_1$ conditions (mean of $80$ trials; change at $t=15$).}\end{figure}
  \begin{figure}[h]\centering\includegraphics[width=\textwidth]{atlas_t29_#1_behavior.pdf}%
    \caption{\textbf{#2} --- psychometric (A,B), chronometric (C), and the cued-quadrant
    attention timecourse (D, the readable summary of the per-frame maps).}\end{figure}
  \clearpage}

\atlast{v11_part2}{v11\_part2 — split-readout cross-talk}{%
  The flagship long-task model and the closest to the paper's attention story. The salience
  stream ($Q{=}X$, $K{=}V{=}[H_1\|H_2]$) feeds the actor; the top-down stream ($Q{=}X$,
  $K{=}V{=}H_2$) feeds the critic. The cued-quadrant timecourse (panel~D) shows the salience
  attention elevated at the cue and the top-down attention rising at the change; the
  reaction time is strongly graded. The cued psychometric sits left of the uncued (a genuine
  cueing benefit).}

\atlast{v11}{v11 — parallel dual-stream}{%
  Two independent streams sharing only the image: salience ($Q{=}X$, $K{=}V{=}H_1$, grounded
  by an $X$ residual) and top-down ($Q{=}X$, $K{=}V{=}H_2$); the heads read both transformer
  outputs. The salience stream carries the cue prior, the top-down stream the change-lock.}

\atlast{v11_part5}{v11\_part5 — split-readout cross-talk (distractor-trained)}{%
  Architecturally identical to \code{v11\_part2} but trained on the cued-\emph{side} task with
  distractors (the uncued side changes on half the trials and must be suppressed). Evaluated
  here on the standard cued task with distractors off, it is among the strongest performers
  (mean cued-hit $0.75$).}

\atlast{v5}{v5 — memory-as-tokens self-attention}{%
  Each encoder layer self-attends over the concatenation $[X\|H_1\|H_2]$ ($300$ tokens). The
  image-key map (panel rows ``image'') tiles the four Gabor quadrants as they appear, and a
  memory hotspot tracks the change; the first encoder layer acts as a cue-driven spatial
  orienter.}

\atlast{v5_part2}{v5\_part2 — cross-attention over memory}{%
  Cross-attention with $Q{=}X$, $K{=}V{=}[X\|H_1\|H_2]$ and two stacked LSTMs read by 1D-conv
  heads. The load-bearing computation is in the memory cross-attention: the memory-key maps
  carry the decision-relevant structure while the image keys provide the spatial substrate.}

\atlast{v8}{v8 — cross-attention with an \texorpdfstring{$H_1$}{H1} residual}{%
  \code{v5\_part2} with the cross-attention residual set to the previous-frame memory $H_1$
  instead of the image, so visual information can reach the decision \emph{only} through the
  attended values. A dedicated patch-reading head emerges and reaction times are fast; the
  image pathway becomes co-dominant with memory.}
